\setauthor{AC}
\section{Projektplanung und Arbeitsweise}
Die Umsetzung des Projekts erfolgte arbeitsteilig innerhalb des Teams. Mein Schwerpunkt lag auf dem Frontend, während mein Kollege die backendseitigen Grundlagen, die Datenhaltung und die serverseitige Logik übernahm. Diese Aufteilung war für das Projekt sinnvoll, da sich sowohl die technischen Anforderungen als auch die Entwicklungsaufgaben deutlich zwischen Benutzeroberfläche und serverseitiger Verarbeitung unterschieden. Gleichzeitig war es notwendig, beide Bereiche regelmäßig aufeinander abzustimmen, damit Datenmodelle, Schnittstellen und Benutzerführung konsistent zusammenarbeiten konnten.

Zu Beginn des Projekts wurde eine To-do-Liste erstellt, um die wichtigsten Funktionen und Arbeitsschritte zu strukturieren. Diese frühe Planung half dabei, den Umfang des Projekts greifbar zu machen und erste Schwerpunkte festzulegen. Im weiteren Verlauf wurde die Arbeitsweise jedoch informeller. Verbesserungen, neue Anforderungen und technische Anpassungen ergaben sich häufig direkt aus der laufenden Umsetzung. Dadurch entstand kein streng linearer Ablauf, sondern ein schrittweiser Entwicklungsprozess, bei dem Planung und Umsetzung eng miteinander verbunden waren.

Die Abstimmung im Team erfolgte unterschiedlich, je nach aktueller Arbeitsphase. Da ein persönlicher Austausch in der Schule häufig möglich war, konnten viele Fragen direkt vor Ort geklärt werden. Ergänzend dazu wurde online per Chat kommuniziert, insbesondere wenn technische Probleme, Schnittstellenfragen oder konkrete Detailentscheidungen zwischen Frontend und Backend abgestimmt werden mussten. Diese Mischung aus direkter und digitaler Zusammenarbeit erwies sich für das Projekt als praktikabel, da sie kurze Rückkopplungen ermöglichte, ohne dass ein komplexes Projektmanagement-Werkzeug erforderlich war.

Auch die Entwicklung im Git-Repository zeigt diesen schrittweisen Projektverlauf. Frühere Commits verweisen zunächst auf grundlegende UI- und Navigationsarbeiten wie die Suchfunktion, das Styling oder Login- und Registeransätze. Darauf aufbauend folgten Phasen, in denen zentrale Funktionen wie Trending-Filme, Abo-Verwaltung und Keycloak-Anbindung umgesetzt wurden. Spätere Commit-Nachrichten wie \glqq Moviesearch überarbeitet und bessere Filterung eingebaut\grqq{}, \glqq Abo Verwaltung hinzugefügt und Bugfixes\grqq{} oder \glqq Keycloak Frontend fertiggestellt\grqq{} machen deutlich, dass die Entwicklung nicht aus isolierten Einzelschritten bestand, sondern aus mehreren aufeinander aufbauenden Verbesserungsphasen.

\section{Technische und fachliche Herausforderungen}
Im Verlauf des Projekts traten mehrere Herausforderungen auf, die sowohl die technische Umsetzung als auch die fachliche Struktur der Anwendung betrafen. Eine zentrale Schwierigkeit bestand für mich darin, die Benutzeroberfläche nicht nur optisch umzusetzen, sondern sinnvoll mit der zugrunde liegenden Logik zu verbinden. Gerade bei der Abo-Verwaltung stellte sich die Frage, welche Informationen für Benutzer tatsächlich relevant sind, wie sie dargestellt werden sollen und an welcher Stelle Interaktion, Übersichtlichkeit und funktionale Tiefe im Gleichgewicht stehen.

Ein weiterer anspruchsvoller Punkt betraf die Kommunikation zwischen Frontend und Backend. Dabei ging es nicht nur darum, Daten technisch abrufen zu können, sondern auch darum, gemeinsam festzulegen, welche Informationen für das Projekt wirklich notwendig sind. Für das Frontend war diese Frage besonders wichtig, da nicht jede backendseitig verfügbare Information automatisch auch einen Mehrwert in der Benutzeroberfläche erzeugt. Die Auswahl relevanter Datenstrukturen, DTOs und Schnittstellen war daher ein wiederkehrender Abstimmungsprozess.

Zusätzlich ergaben sich technische Probleme durch Angular-Versionen und Änderungen innerhalb des Frameworks. Einzelne Implementierungen mussten angepasst oder überarbeitet werden, weil sich Verhaltensweisen, Bibliotheken oder Kompatibilitäten veränderten. Dadurch entstand stellenweise der Eindruck, dass funktionierender Code nicht dauerhaft stabil blieb, sondern nach Framework-Änderungen erneut überprüft werden musste. Diese Erfahrung war insbesondere im Frontend relevant, da moderne Webentwicklung stark von sich rasch weiterentwickelnden Technologien geprägt ist.

Eine weitere Herausforderung lag in der Keycloak-Integration. Zwar konnte die Authentifizierung grundsätzlich erfolgreich eingebunden werden, jedoch war die praktische Verbindung zwischen Login, Token-Verwaltung, geschützten Routen und geschützten Backend-Requests mit mehreren Umsetzungsschritten verbunden. Die Schwierigkeit bestand hier weniger in einer einzelnen Funktion, sondern vielmehr im zuverlässigen Zusammenspiel mehrerer technischer Komponenten. Aus heutiger Sicht war gerade diese Herausforderung besonders lehrreich, weil sie gezeigt hat, wie eng Benutzerführung, Sicherheit und Datenzugriff in einer modernen Webanwendung miteinander verknüpft sind.

Auch die Git-Entwicklung verdeutlicht diese Phase technischer Iteration. Commits wie \glqq durchs Reskripten sind 3 Funktionen kaputt\grqq{}, \glqq es rennt jetzt, paar Funktionen wurde noch nicht implementiert aber großer Teil wurde umgescriptet\grqq{} oder \glqq Filme id hinzugefügt und Abos Logik eingebaut, jedoch fehlt einiges und funktioniert noch nicht\grqq{} zeigen, dass der Projektverlauf nicht rein geradlinig war. Vielmehr mussten funktionierende Teilstände immer wieder überarbeitet, stabilisiert und weiterentwickelt werden. Gerade diese Zwischenphasen waren für den realen Entwicklungsprozess charakteristisch.

\section{Entscheidungsfindung im Projektteam}
Die Entscheidungsfindung im Team erfolgte situationsabhängig. Viele technische und fachliche Fragen wurden gemeinsam diskutiert, insbesondere dann, wenn sie Auswirkungen auf beide Projektteile hatten. Dazu zählten vor allem Fragen der Datenstruktur, der Schnittstellen zwischen Frontend und Backend sowie die Festlegung, welche Funktionen für das Projekt wirklich im Mittelpunkt stehen sollten. Durch diese Diskussionen konnten unterschiedliche Sichtweisen zusammengeführt werden, bevor konkrete Umsetzungen begonnen wurden.

Gleichzeitig arbeitete jeder im eigenen Verantwortungsbereich weitgehend selbstständig. In meinem Fall betraf dies vor allem Entscheidungen zur Gestaltung der Benutzeroberfläche, zur Struktur von Komponenten und zur Einbindung von Authentifizierung und Benutzerinteraktion im Frontend. Obwohl die Umsetzung eigenständig erfolgte, wurde bei wichtigen Fragen regelmäßig Feedback eingeholt. Diese Arbeitsweise war sinnvoll, da sie einerseits Eigenverantwortung ermöglichte und andererseits verhinderte, dass Frontend und Backend inhaltlich auseinanderliefen.

Inhaltlich spielte auch die Abgrenzung des Projektumfangs eine wichtige Rolle. Ein Beispiel dafür war die Frage, welche Inhalte auf der Website tatsächlich angezeigt werden sollten. Im Verlauf der Abstimmung wurde unter anderem diskutiert, ob sich die Anwendung nur auf Filme oder zusätzlich auch auf Serien und Fernsehinhalte beziehen soll. Solche Entscheidungen waren nicht bloß Detailfragen, sondern beeinflussten unmittelbar die Ausrichtung der Anwendung und damit auch die technische Umsetzung. Rückblickend zeigte sich, dass die bewusste Eingrenzung auf Filme wesentlich dazu beitrug, den Projektfokus klar zu halten.

\section{Einbindung von Betreuern und Feedback}
Die Einbindung von Betreuern und Lehrpersonen war im Projekt vor allem dort hilfreich, wo eine Außenperspektive auf Prioritäten, Umfang und inhaltliche Ausrichtung benötigt wurde. Auch wenn nicht jede Projektentscheidung direkt von außen beeinflusst wurde, half das erhaltene Feedback dabei, unklare Punkte schneller zu konkretisieren. Besonders wichtig war dies in Situationen, in denen zwar bereits technische Möglichkeiten vorhanden waren, aber noch nicht eindeutig feststand, welche davon für die Diplomarbeit tatsächlich den größten Nutzen bringen würden.

Das Feedback bezog sich unter anderem auf die Frage, welche Inhalte im Projektkern stehen sollen und welche Funktionen für die Anwendung wirklich relevant sind. Gerade bei der Filmsuche und bei der Auswahl anzuzeigender Inhalte war diese Unterstützung hilfreich, da dadurch der Fokus der Anwendung geschärft wurde. Ebenso wurde deutlich gemacht, dass nicht jede theoretisch denkbare Erweiterung in den Projektumfang aufgenommen werden muss, wenn sie den Kern der Arbeit verwässern würde.

Aus meiner Sicht war diese Form der Einbindung besonders wertvoll, weil sie nicht nur auf technische Details abzielte, sondern vor allem bei der Priorisierung half. Damit wurde der Projektverlauf nicht grundlegend umgelenkt, aber an mehreren Stellen präzisiert. Genau diese punktuelle fachliche Lenkung war für ein Projekt mit mehreren möglichen Entwicklungspfaden hilfreich.

\section{Reaktion auf Fragen und Kritik aus Präsentationen}
Formale Zwischenpräsentationen mit umfangreicher Kritik im klassischen Sinn fanden im Projektverlauf nicht in stark strukturierter Form statt. Rückmeldungen kamen jedoch von Kollegen und aus dem erweiterten schulischen Umfeld, wodurch trotzdem eine wertvolle Außenperspektive entstand. Diese Rückmeldungen bezogen sich häufig auf praktische Fragen der Nutzung und Sicherheit, also genau auf jene Aspekte, die für eine anwendungsnahe Weblösung besonders relevant sind.

Ein Beispiel dafür war die Frage, welche Bereiche durch Keycloak geschützt werden sollten und welche Daten öffentlich zugänglich bleiben können. Solche Hinweise waren wichtig, weil sie halfen, die Grenze zwischen offenen Filmfunktionen und persönlicher Abo-Verwaltung klarer zu ziehen. Auch wenn diese Rückmeldungen nicht in Form formaler Meilenstein-Präsentationen stattfanden, hatten sie dennoch Einfluss auf die Präzisierung einzelner Funktionen.

Insgesamt zeigt sich, dass Kapitel 7 weniger durch starre Projektmanagement-Strukturen geprägt war, sondern stärker durch einen realen, iterativen Entwicklungsprozess. Planung, technische Umsetzung, Abstimmung und Feedback liefen nicht strikt getrennt voneinander ab, sondern beeinflussten sich gegenseitig. Gerade darin liegt ein wesentlicher Erkenntnisgewinn des Projekts: Die Entwicklung einer fachlich sinnvollen Anwendung besteht nicht nur aus dem Schreiben von Code, sondern ebenso aus wiederholtem Abgleichen von Anforderungen, Nutzen, technischer Machbarkeit und Benutzerperspektive.
